% !TEX root = main.tex
\section{Multi-hop reasoning models}\label{appendix:models}

Our \modelshort covers six models with three published works, \ie,  GQE~\cite{hamilton2018embedding}, Q2B~\cite{ren2020query2box} and BetaE~\cite{ren2020beta}, as well as three new methods based on single-hop link prediction models. 
% 
In order to perform logical reasoning in the embedding space, all methods design a projection operator $\Pcal$ and intersection operator $\Ical$. As introduced in ~\Secref{sec:background}, $\Pcal$ represents a mapping from a set of entities (represented by an embedding) to another set of entities (also represented by an embedding) with one relation, \ie, $\Pcal: \mathbb{R}^d\times\Rcal\to\mathbb{R}^d$, assuming the embedding dimension is $d$. The $\Ical$ takes as input multiple embeddings and outputs the embedding that represents the intersected set of entities: $\Ical: \mathbb{R}^d\times\dots\times\mathbb{R}^d\to\mathbb{R}^d$. Different models may have different instantiations of these two operators.

GQE \citep{hamilton2018embedding} is the first method that handles conjunctive queries. GQE embeds a query $q$ to a point in the vector space by iteratively following the computation plan of the query. 
%  
The distance function between the query embedding and the node embedding is the $L_2$ distance: $\texttt{Dist}(\Em{q},\Em{v})=\|\Em{q}-\Em{v}\|_2$. 
% 
Query2box (Q2B)~\cite{ren2020query2box} proposes to embed queries as hyperrectangles with a center embedding and an offset embedding. Q2B simulates set intersection using box intersection. The distance function in Q2B is a modified L1 distance: a weighted sum of the in-box distance and out-box distance.
BetaE~\cite{ren2020beta} further proposes to embed queries as multiple independent Beta distributions so that it can faithfully handle conjunction and negation operation, and thus disjunction by the De Morgan's laws in distribution space with $KL$-divergence as distance. 


Based on the existing methods, we also extend several link prediction models to perform multi-hop reasoning, including RotatE~\cite{sun2019rotate}, DistMult~\cite{yang2014embedding} and ComplEx~\cite{trouillon2016complex}. 

RotatE models a relation by rotation in the complex space, which is naturally compositional. We extend RotatE to RotatE-m for multi-hop reasoning. The relation projection $\Pcal$ is defined as
% 
$\Pcal(\Em{q},\Em{r})=\Em{q}\circ\Em{r}, |\Em{r}[i]|=1$. For the intersection operator, we follow GQE that uses a DeepSets architecture taking multiple complex vectors as input and outputing the intersected result. 
% 
The same distance function in \citet{sun2019rotate} are used, $\texttt{Dist}(\Em{q},\Em{v})=\|\Em{q}-\Em{v}\|$.


Following the design of multihop DistMult proposed in the GQE~\cite{hamilton2018embedding}, we also extend ComplEx to perform multi-hop reasoning. The relation projection is defined as
% 
$\Pcal(\Em{q},\Em{r})=\Em{q}\circ\Em{r}$. And we follow GQE and RotatE-m using the same DeepSets architecture for the intersection operator. The distance function (negative of similarity score) is defined as: $\texttt{Dist}(\Em{q},\Em{v})=-\text{Re}(\langle\Em{q},\overline{\Em{v}}\rangle)$. One key problem of this extension to these semantic matching models is that since the embeddings are not normalized, the similarity score tends to explode especially in the presence of multi-hop relation projection. Thus, we normalize the query embeddings after each step of relation projection or intersection and propose DistMult-m and ComplEx-m. For DistMult-m, we directly perform $L_2$ normalization: $\Em{q}=\frac{\Em{q}}{\|\Em{q}\|_2}$; for ComplEx-m, we normalize the real and imaginary parts separately to the unit hypersphere: 
$\text{Re}(\Em{q})=\frac{\text{Re}(\Em{q})}{\|\text{Re}(\Em{q})\|_2}, \text{Im}(\Em{q})=\frac{\text{Im}(\Em{q})}{\|\text{Im}(\Em{q})\|_2}$. 
